\section{Introduction}
\label{sec:introduction}

Portfolio optimization is a fundamental problem in financial decision-making, wherein investors seek to allocate capital among multiple assets to achieve an optimal balance between expected return and investment risk \cite{markowitz_portfolio_1952}. Since the seminal contribution of Markowitz, traditional portfolio theory has been built on the mean-variance framework, which models risk via the variance of portfolio returns. Despite its mathematical elegance and interpretability, the mean-variance framework rests on strong assumptions that are frequently violated in practice: normally distributed returns, constant market conditions, and fixed investment horizons \cite{campbell_econometrics_1998,demiguel_optimal_2009}. Financial markets are characterised by time-varying dynamics, volatility clustering, regime changes, and heavy-tailed return distributions---features that imply portfolio choices must be continuously recomputed as market conditions evolve \cite{campbell_econometrics_1998}. This necessitates sequential decision-making methods that can handle dynamically changing environments and manage uncertainty in a rigorous manner.

The modelling of multivariate dependencies in financial asset returns presents a fundamental challenge in portfolio optimization. Standard approaches, such as the Gaussian copula or the Dynamic Conditional Correlation (DCC) model of Engle \cite{engle_dynamic_2002}, impose restrictive elliptical dependence structures that may not adequately represent joint tail behaviour \cite{jondeau_copula-garch_2006,rodriguez_measuring_2007}. Vine copulas, introduced by Joe \cite{joe_multivariate_1997} and developed by Bedford and Cooke \cite{bedford_probability_2001,bedford_vinesnew_2002}, overcome this limitation by decomposing a high-dimensional joint distribution into a cascade of bivariate pair-copulas. The present implementation uses a fully dynamic Student-$t$ D-vine: all active tree edges are retained, degrees of freedom are fixed by edge for identifiability, and neural networks update edge correlations causally through time. This provides nonlinear, tail-dependent state information without repeatedly re-estimating a rolling vine or claiming unimplemented dynamic family switching.

So and Yeung \cite{so_vine-copula_2014} extended the vine copula framework to time-varying conditional dependence, demonstrating that dependence among financial assets varies significantly over time. Sahamkhadam and Stephan \cite{sahamkhadam_portfolio_2023} further showed that mixed vine copula models can improve risk reduction across stressed periods. Yet dependence estimation and portfolio control are distinct questions. Most vine-based allocation studies solve a fresh, myopic mean-variance, CVaR, or expected-utility problem at each date. Such optimisers do not directly learn a policy for the path-dependent interaction among current holdings, future rebalancing costs, leverage, and terminal wealth \cite{bernardi_portfolio_2018}.

Reinforcement learning has emerged as a powerful paradigm for sequential decision-making under uncertainty \cite{sutton_reinforcement_2018}. Casting portfolio allocation as a Markov decision process allows an agent to learn a state-contingent policy directly from repeated interactions with a market environment \cite{jang_deep_2023,jiang_deep_2017}. Benhamou et al. \cite{benhamou_bridging_2020} show how deep reinforcement learning can bridge Markowitz-style planning and sequential control. However, financial RL faces two practical difficulties: the available history is short relative to neural-network capacity, and an unconstrained return objective may reward unstable leverage, excessive turnover, or exposure to rare losses. These problems require both a defensible training-data protocol and an economic objective that accounts for downside risk and implementation.

Most portfolio RL states are built from returns, volatilities, and current holdings \cite{jiang_deep_2017,yu_model-based_2020}; the joint conditional distribution is typically implicit. Conversely, a copula generator can expose an agent to more paths, but synthetic data are useful only if their marginal, cross-sectional, tail, and temporal properties are validated. We therefore separate the roles of the dependence model. The neural D-vine (i) generates every episode used in synthetic pre-training, (ii) supplies explicit dynamic dependence features, and (iii) generates conditional scenarios from which the environment computes a policy-visible CVaR statistic and a reward penalty. Historical data ending before the locked final 24 decision periods are used for one fine-tuning pass; no synthetic ``fine-tuning'' sample and no evaluation return enter training.

This paper makes three methodological contributions. First, it integrates a fully dynamic, all-tree neural Student-$t$ D-vine with a recurrent TD3 controller, while keeping the dependence family and edge degrees of freedom fixed so that the claimed dynamics correspond exactly to the implemented correlation forecasts. Second, it formulates a self-financing long--short action map with net exposure one, gross exposure at most 1.5, individual weights in $[-0.20,0.60]$, proportional turnover costs, and explicit short- and cash-borrow charges. Dense reward increments telescope to terminal-wealth CRRA utility, so the objective is genuinely multi-period rather than a sequence of unrelated one-period utilities. Third, it imposes a leakage-resistant two-stage protocol: 1,000 synthetic 24-month episodes are used for pre-training, 61 realised training-prefix trajectories are used once for fine-tuning, and checkpoints must pass numerical, behavioural, and no-holdout gates before evaluation.

The empirical contribution is deliberately evidence-calibrated. The seven-asset adjusted-level panel runs from December 2013 to July 2026, and the final 24 monthly decisions are locked for evaluation. Twenty independently trained policies are combined by averaging target weights before returns and costs are recomputed. Against six causal benchmarks, the ensemble has the lowest observed volatility and highest observed Sharpe ratio over the 22 periods that satisfy the preregistered completeness rule, but it does not maximise terminal wealth or CRRA certainty-equivalent return and its utility effects are not statistically significant. Individual-seed dispersion is material, and much of the ensemble's risk advantage arises because opposing leveraged positions cancel in weight space. A frozen post-holdout explanatory ablation also finds no descriptive advantage over a coarse no-vine TD3 control and almost no increment from historical fine-tuning. We therefore report a mixed result, not a superiority claim.

These findings make causal attribution part of the research contribution. A prospective matched-seed analysis separately removes the direct neural-vine state, the scenario-CVaR observation, the CVaR reward term, synthetic pre-training, recurrence, and historical fine-tuning; it also compares the vine generator with a moving-block bootstrap and TD3 with DDPG, SAC, PPO, and A2C. All variants share the same realised return path, constraints, costs, environment-interaction budget, and ensemble construction. Because the original holdout has already been consumed, these mechanism results are explicitly post-holdout explanatory; fresh confirmation requires frozen non-overlapping walk-forward windows, an external investable panel, or future observations.

The remainder of this paper is organised as follows. Section \ref{sec:methodology} presents the neural D-vine, synthetic-data generator, LSTM-TD3 architecture, state and action design, multi-period objective, training protocol, benchmarks, and causal analysis contract. Section \ref{sec:data} describes the data and sample partition. Section \ref{sec:results} reports the frozen main evaluation, statistical inference, seed robustness, implementation effects, and explanatory ablations. Section \ref{sec:conclusion} concludes and identifies the evidence required for external and confirmatory validation.
